\documentclass[11pt]{article}
\usepackage[utf8]{inputenc}
\usepackage[T1]{fontenc}
\usepackage{amsmath,amssymb,amsthm}
\usepackage{hyperref}
\usepackage{listings}

\title{Anders: Homotopical Library}
\author{Groupoid Infinity}
\date{May 2026}

\begin{document}

\maketitle

\begin{abstract}
This article describes the homotopical library of the Anders proof assistant.
The library provides a comprehensive collection of Higher Inductive Types (HITs),
homotopy groups of spheres, and foundational constructions for synthetic homotopy theory
in a cubical univalent framework.
\end{abstract}

\newpage
\tableofcontents

\newpage
\section{Introduction}

The Anders \footnote{\url{https://github.com/groupoid/anders}} Homotopy Base Library
\footnote{\url{./library/mathematics/homotopy}} implements the core results of synthetic
homotopy theory as presented in the HoTT Book \cite{hottbook} and the Coq/HoTT library \cite{bauer2016}. Leveraging the native
cubical foundations of Anders (CCHM), the library provides direct and computable
implementations of circles, spheres, pushouts, and truncations.

\section{Higher Inductive Types (HITs)}

Higher Inductive Types are a cornerstone of Homotopy Type Theory. In Anders, HITs are implemented
using cubical primitives, which ensure that they are fully computable and satisfy definitional
reduction rules for both point and path constructors. Unlike systems that rely on axiomatic
postulates for HITs, Anders provides a native environment where HIT induction principles
are derived from the underlying cubical structure.

The library includes:
\begin{itemize}
    \item \textbf{Coequalizers}: The quotient of a type by a pair of maps.
    \item \textbf{Hub-Spoke Disc}: A primitive for higher-dimensional glueing.
    \item \textbf{Colimits}: General homotopy colimits over graphs.
    \item \textbf{Pushouts}: $A \sqcup^C B$.
    \item \textbf{Suspensions}: $\Sigma A$.
    \item \textbf{Spheres}: $S^0$, $S^1$, ..., $S^n$.
    \item \textbf{Truncations}: $\| A \|_n$ for various $n$.
    \item \textbf{Eilenberg-MacLane Spaces}: $K(G, n)$.
\end{itemize}

\newpage
\section{Hypercubes}

To support foundational path constructions and complex homotopy proofs, the library includes a module for higher-dimensional paths (\texttt{constcubes.anders}). These are defined recursively as paths between lower-dimensional path primitives:
\begin{itemize}
    \item \textbf{Interval}: A path between two points (1-dimensional).
    \item \textbf{Square}: A path between two paths (2-dimensional).
    \item \textbf{Cube}: A path between two squares (3-dimensional).
    \item \textbf{Tesseract}: A path between two cubes (4-dimensional).
\end{itemize}
These definitions utilize \texttt{PathP} and \texttt{hcomp} to ensure that they are fully computable and can be used to establish the commutativity and associativity of higher loops definitionally.

\section{Coequalizers}

Coequalizers are a foundational construction in the library, representing the quotient of a type $B$ by the equivalence relation generated by two maps $f, g : A \to B$. Formalized in \texttt{coequalizer.anders}, the coequalizer type is denoted:
\begin{equation}
    \text{coeq}(A, B, f, g)
\end{equation}
It includes the constructor \texttt{iota} ($B \to \text{coeq}$) and the path constructor \texttt{resp} which identifies $f(a)$ and $g(a)$ for every $a \in A$. This construction is used to define the circle $S^1$ and other more complex HITs.

\section{Hub-Spoke-Disc Construction}

The Hub-Spoke construction is another foundational primitive in Anders for defining HITs with higher-dimensional boundaries. Detailed in \texttt{hs.anders}, it allows for the creation of a "disc" that fills a sphere $S$:
\begin{equation}
    \text{disc}(S, A)
\end{equation}
This construction provides a center (the base type $A$), a hub (for a map $f : S \to \text{disc}(S, A)$), and spokes (paths between the hub and the image of the map). This is a powerful tool for establishing coherence between higher paths and is used as a building block for truncations and higher groupoids.

\section{Homotopy Colimits}

General homotopy colimits are implemented in \texttt{colim.anders} using a graph-based approach. A diagram is defined by a set of vertices $Vtx$, a set of edges $Edge$, source/target maps $s, t : Edge \to Vtx$, a family of types $D$ over $Vtx$, and actions $act$ over edges. The colimit is then defined as a coequalizer:
\begin{equation}
    \text{colim}(V, E, s, t, D, act) :\equiv \text{coeq}(\Sigma_{e:E} D(s(e)), \Sigma_{v:V} D(v), \sigma, \tau)
\end{equation}
where $\sigma$ and $\tau$ are the source and target maps into the disjoint union of types. This general construction encompasses sequential colimits, pushouts, and other categorical colimits as special cases.

\section{HIT Construction and the Three-HIT Theorem}

A central result in the design of the Anders homotopical library is the ability to construct a wide class of Higher Inductive Types from a small set of foundational primitives. Following the "Three-HIT Theorem"
\cite{ThreeHIT} (and related results in univalent foundations), the library leverages:
\begin{itemize}
    \item \textbf{W-types}: For representing the branching structure of inductive types.
    \item \textbf{Coequalizers}: For identifying points and imposing 1-dimensional path relations.
    \item \textbf{Hub-Spoke Disc}: For filling higher-dimensional spheres and ensuring coherence.
    \item \textbf{Colimits}: For building types through sequential or graph-based processes.
\end{itemize}
By combining these primitives, we can derive complex structures such as $n$-truncations, pushouts, and suspensions without introducing additional axioms. This approach ensures that all HITs in the library share a common, computable foundation and maintain strict adherence to cubical reduction rules.

\newpage
\section{Pushouts and Suspensions}

Pushouts are implemented as the general glueing construction for two maps $f : C \to A$ and $g : C \to B$. Suspensions $\Sigma A$ are then defined as a special case of pushouts:
\begin{equation}
    \Sigma A \equiv \text{pushout}(\mathbf{1}, \mathbf{1}, A, !_A, !_A)
\end{equation}
These constructions form the basis for building higher-dimensional spheres and are fundamental to the calculation of homotopy groups.

\section{Spheres}

The circle $S^1$ is formalized in \texttt{S1.anders} using the coequalizer primitive:
\begin{equation}
    S^1 :\equiv \text{coeq}(\mathbf{1}, \mathbf{1}, \text{id}_{\mathbf{1}}, \text{id}_{\mathbf{1}})
\end{equation}
This provides the base point \texttt{baseS1} and the loop \texttt{loopS1}. Higher spheres $S^n$ are defined by iterating the suspension operation starting from $S^0 \equiv \mathbf{2}$.

\section{Truncations and n-Types}

The library provides n-truncations $\| A \|_n$ for any $n \ge -1$. These are implemented as HITs that "squash" the higher homotopy groups of a type. The truncation operation is used to define:
\begin{itemize}
    \item \textbf{Mere Propositions}: $\| A \|_{-1}$ (also referred to as $hProp$).
    \item \textbf{Sets}: $\| A \|_0$ (also referred to as $hSet$).
    \item \textbf{Surjectivity}: A map $f : A \to B$ is surjective if for all $b \in B$, the fiber is merely inhabited.
\end{itemize}

\section{Loop Spaces and Homotopy Groups}

Higher homotopy groups $\pi_n(A)$ are defined as the truncation of iterated loop spaces:
\begin{equation}
    \pi_n(A, a) :\equiv \| \Omega^n(A, a) \|_0
\end{equation}
The proof of $\pi_1(S^1) \cong \mathbb{Z}$ in \texttt{homotopy.anders} utilizes the encode-decode method over the helix, providing a computable winding number extraction. The library also provides the group structure for these spaces.

\section{Eilenberg-MacLane Spaces $K(G, n)$}

Eilenberg-MacLane spaces $K(G, n)$ are constructed recursively in \texttt{KGn.anders} using the suspension tower and n-truncations:
\begin{equation}
    K(G, n+1) :\equiv \| \Sigma K(G, n) \|_{n+1}
\end{equation}
These types are characterized by having $G$ as their only non-trivial homotopy group at level $n$.

\section{The Hopf Fibration}

The Hopf fibration is formalized in \texttt{hopf.anders} as a bundle $S^1 \to S^3 \to S^2$. The construction leverages univalence to identify the group of isomorphisms of $S^1$ with its loop space, effectively using the rotation map to glue the fibers over the suspension base.

\section{Conclusion}

The Anders Homotopical Library provides a robust foundation for synthetic homotopy theory. By utilizing cubical foundations, the library avoids manual transport along paths, resulting in cleaner proofs and better performance in the type-checker. The library continues to expand with more advanced results in higher categorical structures and colimits.

\begin{thebibliography}{9}
\bibitem{hottbook}
The Univalent Foundations Program.
\textit{Homotopy Type Theory: Univalent Foundations of Mathematics}.
2013.

\bibitem{bauer2016}
Andrej Bauer, Jason Gross, Peter LeFanu Lumsdaine, Michael Shulman, Matthieu Sozeau, and Bas Spitters.
\textit{The HoTT Library: A Comprehensive Resource for Homotopy Type Theory in Coq}.
arXiv:1610.04591 [cs.LO], 2016.

\bibitem{ThreeHIT}
Andrej Bauer and Niels van der Weide.
The Three-HITs Theorem.
DOI:10.4230/LIPIcs.CVIT.2016.23, \url{https://5ht.co/three.pdf}.

\end{thebibliography}
\end{document}
